\begin{table}
\caption{Questions answered by this study. Australian-minus-New Zealand differences are positive when the Australian recorded approval rate is higher.}
\label{tab:discussion-question-answers}
\centering
\small
\begin{tabular*}{\textwidth}{@{\extracolsep{\fill}}p{0.31\textwidth}p{0.10\textwidth}p{0.48\textwidth}@{}}
\toprule
Question & Answer & Explanation \\
\midrule
Is Australia's higher recorded rate explained by origins with higher observed Australian rates accounting for larger shares of its decisions? & No. & With the same average origin distribution, the standardised rates were 88.5\% for Australia and 77.6\% for New Zealand. The 10.9-point difference was larger than the 5.3-point difference under the origin shares observed in the data. \\
\midrule
Are decisions distributed similarly across origins? & No. & Australian decisions were more concentrated. The largest origin accounted for 32.1\% of Australian decisions and 23.8\% of New Zealand decisions. The top five accounted for 73.4\% and 58.2\%, respectively. \\
\midrule
Do the same origins tend to contribute large decision volumes in both destinations? & Yes. & Matched-origin shares had a strong rank association (Spearman $\rho=0.858$), although the share of individual origins differed. \\
\midrule
Does the weighting rule change the comparison? & Yes. & The pooled difference ranged from 5.3 points under the observed destination-specific shares to 16.3 points when every matched origin received equal weight. In 2022, changing the weights also changed the direction. \\
\midrule
Did origin composition produce Australia's higher average rate? & No. & The composition component was $-5.6$ points. The distribution of recorded decisions across origins reduced the pooled Australian-minus-New Zealand difference. \\
\midrule
Was the pattern the same every year? & No. & The rate component calculated under the same origin shares was positive in every year, but its size and the observed decision-weighted gap varied from 2022 to 2025. \\
\midrule
Was the annual pattern caused only by origins entering or leaving the data? & No. & Following the same 118 origins across all four years reproduced the component estimates to one decimal place. \\
\midrule
Did Australia retain a higher average rate after accounting for origin and year? & Yes. & The adjusted difference was 9.8 points (95\% CI 6.6--12.9). Substantial residual heterogeneity remained, so the estimate did not represent every origin and year. \\
\bottomrule
\end{tabular*}
\end{table}
